\documentclass[12pt,a4paper]{article}
\usepackage{amsmath,amssymb}
\usepackage{palatino}
\usepackage[T1]{fontenc}
\usepackage[
  bookmarks=true,
  bookmarksopen=true,
  bookmarksnumbered=true,
  pdftitle={Identity Bedrock: Operationalizing Rompel's Minimal Cryptographic Substrate for Sovereign Identity with Physical Presence Proof},
  pdfauthor={Jean-Francois Bouchard},
  pdfsubject={Sovereign Identity, Cryptography, Physical Presence Proof},
  pdfkeywords={one-way functions, digital signatures, sovereign identity, bone conduction, inaccessible entropy, blockchain identity, post-quantum cryptography},
  colorlinks=true,
  linkcolor=black,
  urlcolor=blue,
  citecolor=black,
  pdfpagelayout=TwoPageRight
]{hyperref}
\usepackage{bookmark}
\bookmarksetup{open,numbered,depth=3}
\usepackage{url}
\usepackage{array}
\usepackage{booktabs}
\usepackage{graphicx}
\usepackage{setspace}
\usepackage{geometry}
\usepackage{amsthm}
\usepackage{float}

\geometry{margin=1.25in}
\setstretch{1.15}

\newcolumntype{P}[1]{>{\raggedright\arraybackslash}p{#1}}

\newtheorem{definition}{Definition}
\newtheorem{theorem}{Theorem}
\newtheorem{remark}{Remark}

%% -------------------------------------------------------
\title{\textbf{Identity Bedrock:}\\[0.4em]
\large Operationalizing Rompel's Minimal Cryptographic Substrate\\
for Sovereign Identity with Physical Presence Proof}

\author{Jean-Fran\c{c}ois Bouchard\\[0.3em]
\small PrivateNess Network (co-founder, CTO)\\
\small Emercoin (community contributor)\\
\small \texttt{[contact via privateness.network]}}

\date{March 2026\quad|\quad Version 13\thanks{%
AI disclosure: The author used AI language model assistance for manuscript
drafting and editing. All technical contributions, system design,
implementation, theoretical framing, and binding insights described herein
originate with the author, drawing on nine years of operational experience
deploying Emercoin-based infrastructure.
Priority established on-chain: Emercoin blocks \#780915 and \#780931.}}

%% -------------------------------------------------------
\begin{document}

\maketitle

\begin{abstract}
Identity Bedrock was not built from theory. It was built from a
requirement: eliminate every institutional key beneath the individual.
Following that requirement to its logical floor---asking what is the
minimum cryptographic primitive that supports public authentication---we
arrived at one-way functions as the irreducible foundation. Rompel's
1990 proof that \textsc{owf}s are necessary and sufficient for
existentially unforgeable signatures under adaptive chosen-message attack
(\textsc{sigeuf-cma}) was found after that convergence. He had already
proved we were standing on bedrock.

The physical presence layer was identified the same way. Examining bone
conduction hardware for a practical deployment, we observed that the
replay channel does not exist: the air-bone impedance mismatch
($T \approx 1.1 \times 10^{-3}$) means the authentication channel cannot
be reconstructed from its output. The connection to Haitner et al.'s
inaccessible entropy framework was found after. The physics preceded the
formalism.

The contribution is the binding: \textsc{owf}-grounded signatures, a
chain-agnostic append-only directory, and a physical presence layer whose
replay resistance follows from acoustic physics rather than computational
hardness assumptions. Each primitive exists in prior literature. The
emergent property---identity with no institutional key beneath the
individual, anchored to Bitcoin via Emercoin \textsc{auxpow}, with
unconditional physical presence proof---belongs to the binding.

This is not a paper describing what could be built. The system has been
running since 2020. Priority is established on-chain.

\medskip
\noindent\textbf{Keywords:} one-way functions, digital signatures,
sovereign identity, bone conduction authentication, inaccessible entropy,
blockchain identity, post-quantum cryptography, physical presence proof.
\end{abstract}

\tableofcontents
\newpage

%%------------------------------------------------------------
\section{Introduction}

\subsection{The Problem: Institutional Capture}

Every deployed digital identity system places institutional keys beneath
the user: certificate authorities, identity providers, key registries,
platform operators. These entities can revoke, deny, surveil, or control
identity without the user's consent.

The threat model motivating this work assumes an adversary with
substantial conventional resources, potential political backing, and full
control of traditional identity infrastructure. The central question is:
can identity exist without institutional keys beneath the individual?

\subsection{The Theoretical Answer: Rompel 1990}

Rompel proved that OWFs are both necessary and sufficient for
\textsc{sigeuf-cma} secure signature schemes~\cite{Rompel90}. This is
the provable minimum---nothing strictly weaker supports public
authentication under standard assumptions.

Our thesis: fixing OWF as the cryptographic floor and adding only a
chain-agnostic append-only directory produces an identity construction
that minimizes institutional capture to zero beneath the individual key.

\subsection{The Binding Insight}

This work follows a pattern established by prior binding innovations in
cryptographic systems: the value lies not in the invention of new
primitives but in recognizing that existing primitives, bound in a
specific way, produce emergent properties none of them exhibits alone.
Each component of Identity Bedrock---OWF-based signatures, append-only
directories, bone conduction biometrics---exists independently in prior
literature. The contribution is the binding and the formal argument for
the security properties that binding produces.

\subsection{Contributions}

\begin{itemize}
  \item \textbf{Theory.} Formalization of a Rompel-grounded,
        chain-agnostic identity system. We prove that removing the OWF
        floor collapses authentication, and that weakening the directory
        trades one consensus for another without improving the
        institutional capture property.

  \item \textbf{Systems.} Architecture: Ed25519 or
        ML-DSA~\cite{FIPS204} identity layer with pluggable directory
        substrate. Implementation: Emercoin NVS (Bitcoin-anchored via
        AuxPoW)~\cite{Emercoin} and NESS Fiber~\cite{NESSFiber}
        (announced publisher keys, append-only BoltDB). Deployed and
        publicly verifiable since 2020.

  \item \textbf{Physical Presence Extension.} Identification of bone
        conduction audio as a physical phenomenon satisfying the
        inaccessible entropy property of Haitner et al.~\cite{Haitner10},
        providing a formal argument for replay resistance that prior
        biometric work demonstrates only
        empirically~\cite{HeadSonic25,HCRAuth24}. Integration as a
        fourth authentication factor and the PBAT token specification.
\end{itemize}

\subsection{Paper Organization}

Section~\ref{sec:background} reviews background and related work.
Section~\ref{sec:threat} defines the threat model.
Section~\ref{sec:architecture} presents the Identity Bedrock architecture.
Section~\ref{sec:implementation} describes the deployed implementation.
Section~\ref{sec:physical} introduces the physical presence extension.
Section~\ref{sec:evaluation} evaluates the system.
Section~\ref{sec:discussion} discusses limitations and future work.

%%------------------------------------------------------------
\section{Background and Related Work}
\label{sec:background}

\subsection{Rompel's Theorem (1990)}

From a complexity-theoretic perspective, OWFs, universal one-way hash
functions (UOWHFs), and \textsc{sigeuf-cma} schemes are equivalent in
strength~\cite{NaorYung89}. Rompel~\cite{Rompel90} and later Haitner et
al.~\cite{Haitner10} showed how to construct UOWHFs from arbitrary OWFs
using inaccessible entropy. A modern exposition is given
in~\cite{KatzKoo05}. The equivalence is:
\[
  \textsc{owf} \Leftrightarrow \textsc{sigeuf-cma}
\]
under standard definitions~\cite{Rompel90,Goldreich01}. To our
knowledge, prior work does not explicitly operationalize this equivalence
as the design floor for a deployed identity system.

\subsection{Inaccessible Entropy}

Haitner et al.~\cite{Haitner10} introduce inaccessible entropy to
simplify and improve Rompel's construction. Informally, a function $F$
has inaccessible entropy if the real entropy of $F^{-1}$---the entropy
remaining in the input after observing the output---significantly exceeds
the accessible entropy: the maximum entropy any polynomial-time adversary
can generate among valid preimages.

\begin{definition}[Inaccessible Entropy, informal~\cite{Haitner10}]
\label{def:inac}
Let $F : \{0,1\}^n \to \{0,1\}^m$. The inaccessible entropy of $F^{-1}$
is the gap
\[
  \Delta = H(X \mid F(X)) - \max_{A \in \mathsf{PPT}}
  H(A(F(X); R) \mid F(X))
\]
where $X$ is uniform on $\{0,1\}^n$, $R$ denotes uniform random coins
for $A$, and $A$ is constrained to output elements of $F^{-1}(F(X))$.
\end{definition}

We observe in Section~\ref{sec:physical} that bone conduction
independently satisfies an analogous property---but enforced by physical
rather than computational means, yielding a strictly stronger guarantee.

\subsection{Blockchain-Based Identity Systems}

\paragraph{Namecoin / EmerDNS.} Name-value storage on blockchain. Gap:
identity keys tied to blockchain wallet keys (not chain-agnostic); miners
act as an implicit certificate authority.

\paragraph{W3C DIDs / Ethereum.} DID methods, smart contracts, ENS. Gap:
DAO governance keys, registry administrators, and provider dependencies
introduce institutional keys above or beneath the user.

\paragraph{Sovrin / Hyperledger Indy.} Permissioned validators, credential
schemas. Gap: steward keys control the network.

\paragraph{Identity Bedrock distinction.} Pure OWF floor. Chain-agnostic
separation of identity cryptography from directory consensus. No
institutional key beneath the user by construction.

\subsection{Bitcoin AuxPoW and Merged Mining}

Auxiliary Proof-of-Work allows Emercoin to be merge-mined with Bitcoin,
inheriting a portion of Bitcoin's security for non-monetary applications.
Khovayko~\cite{hack51} introduces the \emph{hack51} coefficient---defined
as $\mathit{attack\_cost}/\mathit{market\_cap}$---as a numerically
measurable metric for 51\% attack resistance. Under this metric,
Emercoin's triple consensus (PoS $+$ PoW $+$ merged mining) produces a
coefficient significantly exceeding~1.0.

\subsection{Bone Conduction Biometric Authentication}

HeadSonic~\cite{HeadSonic25} proposes a bone conduction earphone
authentication system achieving 96.59\% balanced accuracy across 60
subjects and demonstrating resilience against replay and synthesis
attacks. HCR-Auth~\cite{HCRAuth24} demonstrates reliable authentication
via head contact response. These works establish empirically that bone
conduction authentication resists replay attacks. Neither provides a
formal argument for why. We provide that argument in
Section~\ref{sec:bone-formal}.

\subsection{Skycoin Fiber Consensus}

Skycoin Fiber~\cite{Skycoin} uses announced block publisher keys rather
than economic consensus. Blocks must be signed by announced keys; any
number of hostile nodes up to $N-1$ cannot forge blocks without those
keys. The append-only BoltDB backend makes history rewriting impossible
and any key compromise maximally visible.

\subsection{Summary of Differences}

Identity Bedrock differs from all prior work above in three respects
simultaneously. First, it is the only system explicitly grounded in
Rompel's OWF--signature equivalence as a design floor. Second, the
chain-agnostic architecture decouples identity cryptography from
directory consensus. Third, by identifying bone conduction as satisfying
the inaccessible entropy property, we provide the first formal argument
for why bone conduction biometrics resist replay. The PBAT specification
is the first token format carrying physically-bound biometric session
evidence alongside a cryptographic identity proof in a single signed
structure.

%%------------------------------------------------------------
\section{Threat Model and Requirements}
\label{sec:threat}

\subsection{Adversary Capabilities}

\textbf{Assumed:} Substantial conventional resources and political
backing; full control of traditional identity infrastructure (CAs, IdPs,
KYC registries); ability to outlaw software, seize property, and coerce
individuals; classical (polynomial-time) computational adversary.

\textbf{Not assumed:} Cannot break OWF under standard complexity
assumptions; cannot mount a sustained 51\% attack on Bitcoin (quantified
in Section~\ref{sec:eval-attack}); cannot simultaneously compromise
distributed infrastructure globally; cannot violate acoustic physics
(Section~\ref{sec:bone-formal}).

\textbf{Explicit scope.} This work addresses classical adversaries.
Post-quantum migration to ML-DSA (FIPS~204)~\cite{FIPS204} is discussed
in Section~\ref{sec:limitations}.

\subsection{Design Requirements}

\textbf{R1 --- Mathematical irreducibility.} OWF as cryptographic floor.
No strictly weaker primitive beneath.

\textbf{R2 --- No institutional keys.} Individual key is sole
cryptographic root of trust. No CA, IdP, or registry key beneath the
user.

\textbf{R3 --- Chain-agnostic.} Identity cryptography independent of any
specific directory consensus. Portable across append-only substrates.

\textbf{R4 --- Quantifiable security.} Attack costs explicit and
measurable.

\textbf{R5 --- Operational verification.} Deployed system, not
theoretical proposal. Public verification of all claims.

\textbf{R6 --- Physical presence.} For deployments requiring physical
presence proof, the system must provide it unconditionally---not via
algorithmic liveness detection but via physical channel properties.

%%------------------------------------------------------------
\section{Identity Bedrock Architecture}
\label{sec:architecture}

\subsection{The Two-Layer Design}

\textbf{Layer~1: Identity Cryptography (Chain-Agnostic).}
\[
  \text{Seed}\;(\geq 256\text{ bits})
  \xrightarrow{\mathsf{KeyGen}}
  (sk, pk)
\]
\[
  \text{Identity} = (\mathit{id}, pk),\quad
  \text{Authentication} = \mathsf{Sign}_{sk}(\mathit{challenge})
\]
Properties: OWF-equivalent security (Rompel equivalence); independent of
any blockchain; portable across directory substrates. Default
instantiation: Ed25519~\cite{EdDSA}. Migration target:
ML-DSA-65~\cite{FIPS204}.

\textbf{Layer~2: Directory Substrate (Pluggable).}

\begin{definition}[Append-Only Directory with Finality]
\label{def:dir}
A directory $D \subseteq \mathsf{ID} \times \mathsf{PK}$ is
\emph{append-only with finality} if: (i)~once
$\mathsf{Bind}(\mathit{id}, pk) = 1$ is confirmed, it cannot be set
to~$0$ without an attack of quantifiable cost; and (ii)~the binding is
globally verifiable by any party without trusting an intermediary.
\end{definition}

\textbf{Separation invariant.} Identity uses its chosen signature scheme;
the directory uses whatever consensus fits the deployment's threat model.
The same identity key can move between directories. No coupling exists
between the two layers.

\subsection{Authentication Protocol}

\textbf{Setup.}
\begin{enumerate}
  \item User generates $(sk, pk)$ from entropy source.
  \item User publishes $(\mathit{id}, pk)$ to directory $D$.
  \item Wait for finality (directory-dependent).
  \item Binding is global: $\mathsf{Bind}(\mathit{id}, pk) = 1$.
\end{enumerate}

\textbf{Challenge-Response.}
\begin{enumerate}
  \item $V \to P$: fresh nonce $c \xleftarrow{\$} \{0,1\}^\lambda$.
  \item $P \to V$: $\sigma = \mathsf{Sign}_{sk}(c)$.
  \item $V$ verifies:
    \begin{enumerate}
      \item Fetch $pk$ s.t.\ $\mathsf{Bind}(\mathit{id}, pk) = 1$
            from $D$.
      \item Check $\mathsf{Verify}_{pk}(c, \sigma) = 1$.
      \item Accept iff both hold.
    \end{enumerate}
\end{enumerate}

No state beyond $(\mathit{id}, pk)$ is held in the directory. No smart
contracts. Verification is pure mathematics.

\subsection{Security Proof}

\begin{theorem}[Unforgeable Authentication]
\label{thm:auth}
If the underlying signature scheme is \textsc{sigeuf-cma} secure and $D$
is append-only with finality (Definition~\ref{def:dir}), then Identity
Bedrock provides unforgeable authentication against any classical
adversary.
\end{theorem}

\begin{proof}
Suppose adversary $\mathcal{A}$ wins the authentication game: it produces
valid $(\mathit{id}^*, c^*, \sigma^*)$ without holding $sk$. Construct
forger $\mathcal{B}$ against the signature scheme. $\mathcal{B}$ receives
challenge public key $pk^*$ and simulates the directory by answering
binding queries for $\mathit{id}^*$ with
$pk^*$.\footnote{$\mathcal{B}$ simulates the directory locally; it does
not write to any real ledger.} When $\mathcal{A}$ triggers
authentication, $\mathcal{B}$ forwards the \textsc{sigeuf-cma} game
challenge $c^*$ to $\mathcal{A}$. Since $c^*$ was chosen fresh and never
signed by the legitimate holder of $sk$, $\mathcal{A}$'s forgery
$\sigma^*$ constitutes an existential forgery. By security of the
scheme:
\[
  \Pr[\mathcal{A}\text{ wins}] \leq \mathsf{negl}(\lambda).
\]
Directory attacks shift cost to the directory layer, quantified in
Section~\ref{sec:eval-attack}.
\end{proof}

\subsection{Chain-Agnostic Property and Migration}

The same $(sk, pk)$ can be published to Emercoin NVS or any append-only
ledger satisfying Definition~\ref{def:dir}. Migration is authenticated:
the user signs a migration statement with $sk$, publishes to $D'$, and
verifiers accept after finality. The old directory retains the historical
record.

%%------------------------------------------------------------
\section{Implementation: Emercoin + NESS Fiber}
\label{sec:implementation}

\begin{figure}[H]
\centering
\includegraphics[width=0.50\textwidth]{NESS-hexagon.jpg}
\caption{The Identity Bedrock trust stack. NESS (top) anchors to
Emercoin and Skywire (middle), which anchor to Bitcoin via AuxPoW
(bottom). Each layer derives security from the layer beneath it without
introducing institutional keys.}
\label{fig:stack}
\end{figure}

The architecture builds on the PrivateNess Network design published in
October~2020~\cite{Litepaper20}, which established the dual-blind
Emercoin $+$ Skywire architecture, 127.0.0.1 communication restriction,
UHE-PRNG entropy at 1536~bits of internal state, and the full OSI-stack
decentralization approach. Identity Bedrock formalizes the cryptographic
foundations of that deployed system.

\subsection{Directory: Emercoin NVS}

Emercoin NVS provides a name-value store on a hybrid PoS $+$ AuxPoW
blockchain~\cite{Emercoin}. Consensus: 85\% Proof-of-Stake $+$ 15\%
Bitcoin AuxPoW. Finality: AuxPoW confirmation $\approx$10~minutes;
PoS-only blocks arrive faster but carry probabilistic finality only.

Identity bindings are anchored using the WORM (World Object MaRkup
language) schema~\cite{WORM}, a decentralized XML-based search
infrastructure designed by Aleksej Sokolov and deployed within the
PrivateNess ecosystem:

\begin{verbatim}
<worm>
  <user type="ness" key="<pubkey_hex>" nonce="<nonce>">
    <keys>
      <key hash="<pubkey_hash>" algo="ed25519">
        <pubkey_hex>
      </key>
    </keys>
    <node type="ness" url="<skywire_pubkey>"
          nonce="<nonce>" tags="recorder,voice"/>
  </user>
</worm>
\end{verbatim}

The \texttt{<keys>} list natively supports multiple algorithms, enabling
ML-DSA migration without schema changes.

\paragraph{Separation of concerns.} The Emercoin wallet key controls NVS
record updates. The identity key authenticates the identity itself.
Directory control is strictly separated from identity authentication.

\subsection{Infrastructure: NESS Fiber}

NESS Fiber uses Skycoin-type Fiber consensus~\cite{Skycoin} with
announced publisher keys published in Emercoin NVS under
\texttt{dpo/PrivateNESS.Network}~\cite{NESSFiber}, creating a verifiable
trust chain: NESS Fiber $\to$ Emercoin NVS $\to$ Bitcoin AuxPoW. This
anchoring is trustless: any verifier can confirm the binding
independently at the Emercoin block explorer without relying on any
intermediary.

Key compromise is maximally detectable: a stolen publisher key cannot
rewrite past blocks (append-only BoltDB makes history rewriting
impossible). The only action available is a reset to genesis block~0,
immediately visible to all participants.

\subsection{Canonical Self-Spend Proof}

The canonical self-spend proof binds a chain address to the Ed25519
identity:
\begin{enumerate}
  \item Create unsigned self-spend transaction $A \to A$
  \item Sign with the secp256k1 chain key
  \item Create proof $\{\mathit{network}, \mathit{address},
        \mathit{policy}, \mathit{txid}, \mathit{tx\_hex}\}$
  \item Anchor proof in Emercoin NVS WORM record
  \item WORM record controlled by the Ed25519 bedrock key
\end{enumerate}

\textbf{Deployed instance.} Anchored in Emercoin block~\#445669 (PoS
block, November~1, 2020, 22:16:45~UTC). The Bitcoin AuxPoW security
moat is distributed across the Emercoin chain at the chain level; the
protection is structural, not block-specific. Verifiable at
\url{https://explorer.emercoin.com/dpo/PrivateNESS.Network}.

\subsection{Key Generation and Entropy Provisioning}

Identity keys are derived from a seed produced by Gibson's
UHEPRNG~\cite{UHEPRNG}. Entropy sources mixed in include: multiple
1024-byte blocks from \texttt{os.urandom} ($\geq$1536~bits total); RC4OK
output from Emercoin core RNG; CPU timing jitter; and hardware
identifiers (eMMC CID, MAC address, machine-id). Security rests on the
standard $\approx$128-bit classical security floor of the chosen scheme.

%%------------------------------------------------------------
\section{Physical Presence Extension}
\label{sec:physical}

\subsection{Motivation}

The two-layer Identity Bedrock system addresses forgery (OWF layer) and
directory rewrite (AuxPoW layer). It does not address physical presence:
a device with a stolen key can impersonate the legitimate owner. For
deployments where physical presence is a security requirement, an
additional unconditionally guaranteed factor is needed.

\subsection{The Channel Symmetry Problem}

Every voice authentication system built on air-conducted audio shares a
fundamental vulnerability: the authentication and replay channels are
physically identical. A microphone captures air pressure waves. A speaker
produces air pressure waves. This is not a software deficiency---it is a
property of air-conducted sound that no liveness detection algorithm can
fully overcome.

\subsection{Bone Conduction Breaks Channel Symmetry}
\label{sec:bone-formal}

Bone conduction transmits mechanical vibration through skull bone and
soft tissue to the transducer. Replay requires producing the same
mechanical vibration pattern at the transducer---requiring the same skull
geometry, bone density, tissue distribution, and coupling characteristics
as the enrolled individual.

The physical barrier to replay is the mechanical impedance mismatch at
the air-bone interface. The specific acoustic impedance of air is
$Z_\mathrm{air} \approx 415~\mathrm{Pa{\cdot}s/m}$; cortical bone is
$Z_\mathrm{bone} \approx 1.5 \times
10^6~\mathrm{Pa{\cdot}s/m}$~\cite{Kinsler00}. The pressure transmission
coefficient is:
\begin{equation}
  T = \frac{4 Z_\mathrm{air} Z_\mathrm{bone}}
           {(Z_\mathrm{air} + Z_\mathrm{bone})^2}
    = \frac{4 \times 415 \times 1.5 \times 10^6}
           {(415 + 1.5 \times 10^6)^2}
    \approx 1.1 \times 10^{-3}
  \label{eq:transmission}
\end{equation}

Approximately 0.11\% of acoustic energy from an air-conducted source
transmits into bone. The spatial coupling characteristics---which depend
on the geometry of the skull in contact with the transducer---are
entirely absent in any replay attempt.

\begin{remark}
\label{rem:physics}
Equation~\eqref{eq:transmission} is a consequence of acoustic physics,
not a hardness assumption. It holds unconditionally, independent of any
computational model.
\end{remark}

\subsection{Bone Conduction as Physical Inaccessible Entropy}

\begin{remark}[Physical instantiation of inaccessible entropy]
\label{rem:inacc}
The bone conduction channel produces a distribution---the mechanical
resonance signature of the enrolled skull---whose entropy is inaccessible
to any adversary not in physical control of the enrolled body in contact
with the enrolled transducer. The inaccessibility is enforced by
Equation~\eqref{eq:transmission}, not by any computational hardness
assumption. Where Definition~\ref{def:inac} requires inaccessibility
against \textsf{PPT} adversaries (a computational bound), bone conduction
provides inaccessibility against all adversaries lacking physical control
(an unconditional bound)---a strictly stronger guarantee under any
computational model.
\end{remark}

\begin{remark}[Scope of Remark~\ref{rem:inacc}]
Remark~\ref{rem:inacc} is an observation connecting a physical phenomenon
to the inaccessible entropy framework of Haitner et
al.~\cite{Haitner10}. A formal reduction demonstrating that bone
conduction channel entropy satisfies Definition~\ref{def:inac} under the
standard model remains for future work.
\end{remark}

\subsection{The Four-Factor Authentication Stack}

\begin{table}[H]
\caption{Four-factor authentication stack}
\label{tab:stack}
\small\renewcommand{\arraystretch}{1.3}
\begin{tabular}{P{2.8cm}P{2.8cm}P{2.8cm}P{2.8cm}}
\toprule
\textbf{Factor} & \textbf{Mechanism} & \textbf{Attack addressed}
  & \textbf{Guarantee}\\
\midrule
Ed25519 / OWF    & Rompel floor           & Forgery
  & Computational\\
EmerNVS / AuxPoW & Bitcoin-anchored dir.  & Dir.\ rewrite
  & Economic\\
Bone conduction  & Physical coupling      & Replay
  & Unconditional\\
Voiceprint+PBAT  & Session embed.\ signed & Impersonation
  & Biometric+crypto\\
\bottomrule
\end{tabular}
\end{table}

\subsection{The Physically-Bound Authentication Token (PBAT)}

\begin{definition}[PBAT]
\label{def:pbat}
A \emph{Physically-Bound Authentication Token} is a signed data
structure containing: an identity public key; a directory binding proof;
a biometric session embedding captured via a physically inaccessible
channel (Definition~\ref{def:phys}); a single-use nonce; a timestamp;
and a digital signature over all fields by the identity private key.
\end{definition}

\begin{verbatim}
{
  "identity_pubkey":   "<Ed25519 or ML-DSA pubkey hex>",
  "identity_algo":     "ed25519 | mldsa65",
  "nvs_record":        "<EmerNVS WORM txid>",
  "voiceprint_hash":   "<SHA256 or SHA3-256 of enrolled embedding>",
  "session_embedding": "<current session voice embedding>",
  "similarity_score":  0.94,
  "capture_channel":   "bone_conduction_msbc",
  "headset_mac":       "<BT MAC of enrolled headset>",
  "timestamp":         "<UTC ISO8601>",
  "nonce":             "<random 256-bit>",
  "expires_seconds":   300,
  "signature":         "<Sign_sk(all above fields)>"
}
\end{verbatim}

A PBAT cannot be replayed: the nonce is single-use, the timestamp is
signed, and the session embedding is captured live at this authentication
event.

\textbf{Physical presence doctrine.} The device authenticates the body,
not the credential. No network path, no software override, and no root
privilege substitutes for physical presence of the enrolled body wearing
the enrolled transducer.

\textbf{Quantum note.} The physical presence layer is unconditionally
quantum-resistant. Post-quantum migration is scoped to the cryptographic
binding layer only.

%%------------------------------------------------------------
\section{Evaluation}
\label{sec:evaluation}

\subsection{Attack Cost Analysis}
\label{sec:eval-attack}

\begin{table}[H]
\caption{Attack cost analysis}
\label{tab:attacks}
\small\renewcommand{\arraystretch}{1.3}
\begin{tabular}{P{2.5cm}P{3cm}P{2.5cm}P{2.5cm}}
\toprule
\textbf{Attack vector} & \textbf{Requirements} & \textbf{Estimated
  cost} & \textbf{Result}\\
\midrule
Break signature scheme &
  Violate OWF assumption &
  Infeasible (classical) &
  Authentication collapses\\
Rewrite Emercoin NVS &
  PoS majority $+$ PoW $+$ AuxPoW 51\% simultaneously &
  $\approx$5$\times$ EMC mkt.\ cap~\cite{hack51}; Bitcoin
  component $\sim$\$5--20B~\cite{CoinMetrics24} &
  Directory rewrite\\
Compromise NESS Fiber &
  Steal publisher keys &
  High; DoS only without keys &
  Cannot alter identity bindings\\
Replay bone conduction &
  Reproduce mechanical skull signature via air &
  Physically infeasible ($T\approx10^{-3}$,
  Eq.~\eqref{eq:transmission}) &
  Biometric layer holds\\
Coerce enrolled user &
  Physical &
  Variable &
  Outside cryptographic scope\\
\bottomrule
\end{tabular}
\end{table}

Attack cost figures depend on current market conditions and should be
treated as order-of-magnitude illustrations. The structural property
(hack51 coefficient $\gg$1.0) is what matters~\cite{hack51}.

\subsection{Comparison to Related Systems}

\begin{table}[H]
\caption{Comparison of identity systems}
\label{tab:comparison}
\small\renewcommand{\arraystretch}{1.3}
\begin{tabular}{P{2.8cm}P{2.5cm}P{1.8cm}P{2cm}P{2.5cm}}
\toprule
\textbf{System} & \textbf{Identity floor} & \textbf{No inst.\
  keys} & \textbf{Chain-agnostic} & \textbf{Physical presence}\\
\midrule
Identity Bedrock    & Ed25519/Rompel & \checkmark & \checkmark
  & Physical (bone)\\
X.509 PKI           & RSA/ECDSA      & $\times$   & $\times$
  & None\\
Namecoin            & ECDSA          & $\times$   & $\times$
  & None\\
Ethereum DIDs       & Varies         & $\times$   & $\times$
  & None\\
Sovrin              & Ed25519        & $\times$   & $\times$
  & None\\
FIDO2               & ECDSA          & Partial    & $\times$
  & HW token\\
HeadSonic~\cite{HeadSonic25} & None & N/A        & N/A
  & Physical (bone)\\
\bottomrule
\end{tabular}
\end{table}

\subsection{Performance}

\begin{itemize}
  \item Emercoin NVS finality: $\approx$10~minutes (AuxPoW
        confirmation)
  \item Challenge-response: milliseconds (no on-chain transaction at
        verification time)
  \item Voiceprint verification
        (resemblyzer,\footnote{\url{https://github.com/resemble-ai/resemblyzer}}
        RK3528 class hardware): $<$1~second
  \item PBAT generation and signing: $<$10~ms
\end{itemize}

\subsection{Reference Implementation: Minimum Hardware Floor}
\label{sec:refimpl}

A design goal of Identity Bedrock is accessibility: sovereign identity
must not require expensive or specialized hardware.

\begin{table}[H]
\caption{Minimum hardware floor for full PBAT deployment}
\label{tab:hardware}
\small\renewcommand{\arraystretch}{1.3}
\begin{tabular}{P{3.2cm}P{4.2cm}P{2.2cm}}
\toprule
\textbf{Component} & \textbf{Reference device} & \textbf{Cost (CAD)}\\
\midrule
Compute node &
  X88 Pro~13 (RK3528 SoC), repurposed Android TV box, Armbian Linux &
  $\approx$\$38\\
Bone conduction headset &
  ERAZER XF22, Bluetooth~5.4, open-ear skull-contact transducer &
  \$9.69\\
Directory anchor &
  Emercoin NVS record (one-time write) &
  $<$\$1\\
Software &
  Armbian, resemblyzer, Identity Bedrock stack &
  \$0\\
\midrule
\textbf{Total} & & $\mathbf{\approx\$49}$\\
\bottomrule
\end{tabular}
\end{table}

\textbf{PBAT eligibility criterion.} Any bone conduction headset is
PBAT-eligible if and only if: (1)~its transducer couples mechanically to
skull bone (open-ear form factor, not in-ear); (2)~HFP negotiation
succeeds at 16~kHz on the target host (\texttt{arecord} at
\texttt{-r~16000} completes without error); (3)~resemblyzer produces
consistent cosine similarity $\geq 0.85$ across three independent
enrollment sessions from the same individual.

\textbf{Political note.} Sovereign identity with cryptographic grounding,
Bitcoin-anchored directory, and unconditional physical presence proof for
under \$49~CAD total hardware cost is a direct answer to the claim that
privacy-preserving infrastructure is only accessible to the technically
sophisticated or well-resourced. The PrivateNess Network has operated on
this principle since 2020~\cite{Litepaper20}.

%%------------------------------------------------------------
\section{Discussion}
\label{sec:discussion}

\subsection{Limitations}
\label{sec:limitations}

\textbf{Out of scope.} Key compromise via malware or physical theft;
social engineering; UI-layer attacks.

\textbf{Quantum adversary.} Ed25519 is vulnerable to Shor's algorithm.
Migration path: ML-DSA-65 (FIPS~204~\cite{FIPS204}).
\begin{itemize}
  \item Replace Ed25519 with ML-DSA-65 (pubkey 1952~bytes vs.\ 32;
        accommodated by EmerNVS value field)
  \item Replace SHA-256 voiceprint hash with SHA3-256 in PBAT
  \item WORM \texttt{<keys>} schema natively supports multiple
        algorithms; migration adds \texttt{<key algo="mldsa65">}
        without modifying existing records
  \item Physical presence layer, PBAT structure, and directory anchoring
        unchanged
  \item OpenSSL~3.5 (April~2025) includes native ML-DSA support in
        mainline Armbian Linux
\end{itemize}

\textbf{Formal reduction for bone conduction.} Remark~\ref{rem:inacc}
is an observation, not a completed theorem. A formal reduction under the
standard model remains for future work.

\textbf{Attack cost precision.} Absolute figures depend on current EMC
market conditions and should be treated as order-of-magnitude
illustrations.

\textbf{Coercion.} Residual attack surface: coercion, enrollment
compromise, and post-quantum cryptanalysis---each explicitly bounded.

\subsection{On the Nature of the Contribution}

This work follows the pattern of binding innovations in applied
cryptography: taking existing primitives and identifying a binding that
produces emergent properties none of them exhibits alone. The emergent
property is an identity system with no institutional key beneath the
individual, unconditional physical presence proof, and a formal
cryptographic floor, deployed and publicly verifiable since
2020~\cite{NESSFiber}.

\subsection{Future Work}

\begin{itemize}
  \item Formal reduction: prove bone conduction channel entropy satisfies
        Definition~\ref{def:inac} under the standard model
  \item Full PBAT security proof under the standard model
  \item Empirical verification of PBAT eligibility criteria
        (Section~\ref{sec:refimpl}) on the ERAZER XF22 reference device:
        mSBC negotiation confirmation and resemblyzer consistency across
        enrollment sessions
  \item Threshold signatures ($k$-of-$n$ ML-DSA) for social recovery
        without institutional escrow
  \item User studies on bone conduction enrollment UX and false rejection
        rate across hardware price points
  \item Dedicated sovereign identity node hardware specification,
        community-auditable, derived from the RK3528 prototype lineage
\end{itemize}

%%------------------------------------------------------------
\section*{Conclusion}

We have presented Identity Bedrock, a deployed identity system grounded
in Rompel's OWF--signature equivalence---arrived at not from reading the
theorem first, but from following the requirement of zero institutional
capture to its logical floor. The system achieves no institutional keys
beneath the individual through a chain-agnostic architecture separating
identity cryptography from directory consensus.

We have further identified bone conduction audio as a physical phenomenon
satisfying the inaccessible entropy property of Haitner et
al.~\cite{Haitner10} (Remark~\ref{rem:inacc}), providing the physical
argument for replay resistance that prior biometric work demonstrates
empirically. Integrated as a fourth authentication factor, this produces
the four-factor system of Table~\ref{tab:stack}.

The canonical on-chain identity record at~\cite{NESSFiber} constitutes
operational proof: the genesis address, blockchain public key, and
genesis signature are anchored in Emercoin NVS via AuxPoW, transitively
anchored to Bitcoin, and independently auditable by any party without
trusting an intermediary. This is not a paper describing what could be
built. It is a paper describing what has been running since 2020.

%%------------------------------------------------------------
\section*{Acknowledgments}

The author thanks Oleg Khovayko for foundational Emercoin infrastructure
work and the hack51 attack cost framework; Aleksej Sokolov (Alex014) for
designing the WORM format and implementing the PrivateNess toolchain
(privatenesstools, pyuheprng) from the author's specifications; Brandon
``Synth'' Smietana and Moses Narrow (moe\_narrow / 0pcom), Skywire Linux
packaging lead, for the Skywire mesh transport and Fiber consensus
infrastructure on which NESS is built; and the PrivateNess Network
community, whose collective work since the project's public founding on
Bitcointalk on September~1, 2020~\cite{BTT20} has made this deployed
system possible.

%%------------------------------------------------------------
\begin{thebibliography}{99}

\bibitem{Rompel90}
J.~Rompel.
One-way functions are necessary and sufficient for secure signatures.
In \textit{Proc.\ 22nd Annual ACM STOC}, pages 387--394. ACM, 1990.

\bibitem{Haitner10}
I.~Haitner, T.~Holenstein, O.~Reingold, S.~Vadhan, and H.~Wee.
Universal one-way hash functions via inaccessible entropy.
In \textit{EUROCRYPT~2010}, LNCS vol.~6110, pages 616--637.
Springer, 2010.

\bibitem{EdDSA}
S.~Josefsson and I.~Liusvaara.
Edwards-curve digital signature algorithm (EdDSA).
RFC~8032, IETF, January~2017.
\url{https://www.rfc-editor.org/rfc/rfc8032}

\bibitem{UHEPRNG}
S.~Gibson.
Ultra-high entropy pseudo-random number generator (UHEPRNG).
Gibson Research Corporation.
\url{https://www.grc.com/otg/uheprng.htm}

\bibitem{Emercoin}
Emercoin Development Team.
Emercoin: hybrid blockchain with merged mining.
\url{https://emercoin.com}, 2013--2026.

\bibitem{hack51}
O.~Khovayko.
Estimating the cost of a 51\% attack on Emercoin: the hack51
coefficient [In Russian].
Dzen / Yandex, 2021.
\url{https://dzen.ru/a/YUotZkN9nAeElPEQ}

\bibitem{CoinMetrics24}
L.~Nuzzi et~al.
Breaking BFT: quantifying the cost to attack Bitcoin and Ethereum.
Coin Metrics, 2024.

\bibitem{Skycoin}
B.~Smietana (``Synth'') et~al.
Skycoin: Skywire and Fiber consensus.
Skycoin Development Team, 2014--2026.
\url{https://www.skycoin.com}
Linux packaging: M.~Narrow (moe\_narrow / 0pcom).

\bibitem{WORM}
A.~Sokolov (Alex014) / NESS-Network.
WORM --- World Object MaRkup language.
\url{https://github.com/ness-network/worm}

\bibitem{NaorYung89}
M.~Naor and M.~Yung.
Universal one-way hash functions and their cryptographic applications.
In \textit{Proc.\ 21st Annual ACM STOC}, pages 33--43. ACM, 1989.

\bibitem{KatzKoo05}
J.~Katz and C.-Y.~Koo.
On constructing universal one-way hash functions from arbitrary
one-way functions.
Cryptology ePrint Archive, Report 2005/328, 2005.
\url{https://eprint.iacr.org/2005/328}

\bibitem{FIPS204}
National Institute of Standards and Technology.
Module-lattice-based digital signature standard (ML-DSA).
FIPS~204, August~2024.
\url{https://doi.org/10.6028/NIST.FIPS.204}

\bibitem{HeadSonic25}
Z.~He, J.~Chen, K.~He, Y.~Gu, Q.~Deng, and Z.~Zhang.
HeadSonic: usable bone conduction earphone authentication via
head-conducted sounds.
\textit{IEEE Trans.\ Mobile Comput.}, 24(9):7914--7928, 2025.
\url{https://doi.org/10.1109/TMC.2025.3551272}

\bibitem{HCRAuth24}
Z.~He, J.~Chen, C.~Wu, K.~He, R.~Du, J.~Jia, Y.~Gu, and X.~Sun.
HCR-Auth: reliable bone conduction earphone authentication with head
contact response.
\textit{Proc.\ ACM IMWUT}, 8(4):1--27, 2024.
\url{https://doi.org/10.1145/3699780}

\bibitem{Kinsler00}
L.~E.~Kinsler, A.~R.~Frey, A.~B.~Coppens, and J.~V.~Sanders.
\textit{Fundamentals of Acoustics}, 4th~ed.
Wiley, 2000. ISBN: 978-0-471-84789-2.

\bibitem{Goldreich01}
O.~Goldreich.
\textit{Foundations of Cryptography, Vol.~1: Basic Tools}.
Cambridge University Press, 2001.

\bibitem{BTT20}
PrivateNESS Network.
PrivateNess Network (NESS) Privacy CX Fiber Coin.
Bitcointalk, September~1, 2020.
\url{https://bitcointalk.org/index.php?topic=5272741.0}

\bibitem{Litepaper20}
J.-F.~Bouchard et~al.
PrivateNess Network LitePaper.
PrivateNess Network, October~9, 2020.
\url{https://privateness.network/?p=255}

\bibitem{NESSFiber}
PrivateNESS Network.
PrivateNESS.Network --- canonical on-chain identity record.
Emercoin NVS DPO, block~\#445669, November~1, 2020.
\url{https://explorer.emercoin.com/dpo/PrivateNESS.Network}

\end{thebibliography}

%%------------------------------------------------------------
\appendix

\section{Formal Definitions}

\begin{definition}[One-Way Function]
A function $f : \{0,1\}^* \to \{0,1\}^*$ is a \emph{one-way function}
if it is computable in polynomial time and for all \textsf{PPT}
adversaries $\mathcal{A}$:
\[
  \Pr\!\left[\mathcal{A}(f(x)) \in f^{-1}(f(x))\right]
  \leq \mathsf{negl}(\lambda)
\]
where $x \xleftarrow{\$} \{0,1\}^\lambda$.
\end{definition}

\begin{definition}[\textsc{sigeuf-cma} Security]
A signature scheme $(\mathsf{Gen}, \mathsf{Sign}, \mathsf{Verify})$ is
\emph{existentially unforgeable under adaptive chosen-message attack} if
for all \textsf{PPT} adversaries $\mathcal{A}$ with access to a signing
oracle:
\[
  \Pr\!\left[\mathsf{Verify}(pk, m^*, \sigma^*) = 1
  \;\wedge\; m^* \notin Q\right] \leq \mathsf{negl}(\lambda)
\]
where $(pk, sk) \leftarrow \mathsf{Gen}(1^\lambda)$ and $Q$ is the set
of queried messages.
\end{definition}

\begin{definition}[Physically Inaccessible Biometric Channel]
\label{def:phys}
A biometric capture channel is \emph{physically inaccessible} if the
entropy produced by the channel is inaccessible to any adversary without
physical control of the enrolled subject in contact with the enrolled
transducer, by physical properties of the channel rather than by
computational hardness assumptions.
\end{definition}

\section{WORM Identity Record Schema}

The WORM \texttt{<keys>} list supports multiple algorithms, enabling
post-quantum migration without schema changes:

\begin{verbatim}
<worm>
  <user type="ness" key="<primary_pubkey_hex>"
        nonce="<random_nonce>">
    <keys>
      <key hash="<pubkey_hash>" algo="ed25519">
        <pubkey_hex>
      </key>
      <key hash="<mldsa_hash>" algo="mldsa65">
        <mldsa_pubkey_hex>
      </key>
    </keys>
    <node type="ness" url="<skywire_pubkey>"
          nonce="<skywire_nonce>" tags="recorder,voice"/>
  </user>
</worm>
\end{verbatim}

\section{Deployment Verification}

\begin{itemize}
  \item On-chain identity record (EmerNVS DPO):\\
        \url{https://explorer.emercoin.com/dpo/PrivateNESS.Network}\\
        Block~\#445669, November~1, 2020, 22:16:45~UTC.
  \item Bitcointalk (September~1, 2020):\\
        \url{https://bitcointalk.org/index.php?topic=5272741.0}
  \item Emercoin explorer: \url{https://explorer.emercoin.com}
  \item PrivateNess Network: \url{https://privateness.network}
  \item PrivateNess tools:\\
        \url{https://github.com/ness-network/privatenesstools}
  \item pyuheprng:\\
        \url{https://github.com/ness-network/pyuheprng}
  \item WORM: \url{https://github.com/ness-network/worm}
  \item Block explorers:\\
        \url{https://ness-explorer.magnetosphere.net}\\
        \url{https://explorer.privateness.network}
\end{itemize}

\end{document}
